\section{An exact series and certified value for adjacent CUE agreement}

\begin{theorem}[NN-A017: adjacent agreement from integer gap determinants]
Let $K_s$ be the sine-kernel operator on an interval of length $s$,
and put $E(s)=\det(I-K_s)$, the empty-interval probability for the
unit-intensity stationary sine process. The limiting adjacent-label
agreement probability in NN-A003 is
\[
 A_{\rm adj}=1-2\sum_{n=1}^\infty(-1)^{n+1}E(n).
                                                        \tag{A017.1}
\]
The series is absolutely convergent and its alternating partial sums
give certified two-sided bounds. In the specified model, for every
fixed exterior radius $0\le r<1$ and offset,
\[
 0.66654660183721072003
 < A_{\rm adj} <
 0.66654660183721072008.                                \tag{A017.2}
\]
The interval is certified using degree-$100$ Taylor features and
Arb ball arithmetic at $160$ decimal digits; no quadrature or
uncertified extrapolation is used. It is a CUE-model value, not a
proved proportion for arithmetic zero labels.
\end{theorem}

\begin{proof}
Let $D_1$ be the first positive point under the reduced Palm sine law.
The endpoint derivative of a void probability gives
\[
 -E'(s)=\mathbb P^0(D_1>s),\qquad E''(s)=f_1(s),
                                                        \tag{A017.3}
\]
where the second equality may first be read distributionally.
Indeed, enlarging $(0,s)$ by $ds$ loses precisely configurations
with a point in that added strip and no point in the old interval;
the Campbell formula gives the reduced Palm void probability at
the right endpoint. Translation and reflection symmetry identify
it with the probability on the right side of (A017.3).
The chance of two added points is $O(ds^2)$, since the two-point
intensity is bounded by one. Formula (A016.6) supplies its density;
the two-point Palm kernel varies continuously in trace norm after
rescaling each interval to a fixed unit interval, on every compact
set of $s>0$, so that formula also gives continuity. The distributional
form alone is enough below. In particular $E(0)=1$ and $E'(0)=-1$.

Write $g(s)=d_2(1-s)$ for $s\ge0$. It is the triangular function
with $g(0)=1$, right derivative $g'(0)=-1$, and distributional
second derivative
\[
 g''=2\sum_{n\ge1}(-1)^{n+1}\delta_n.
\]
All eigenvalues of $K_s$ lie in $[0,1]$ and their sum is $s$;
therefore $0\le E(s)\le e^{-s}$. Also $E'(s)\to0$ because
$D_1$ is finite almost surely. Twice integrating by parts, first on
a finite interval and then passing to infinity, gives
\[
 \mathbb E d_2(1-D_1)
 =\int_0^\infty g(s)E''(s)ds
 =2\sum_{n\ge1}(-1)^{n+1}E(n).
\]
The two terms at zero cancel: $-g(0)E'(0)=1$ and
$g'(0)E(0)=-1$. The boundary terms at infinity vanish; the series
converges absolutely by $E(n)\le e^{-n}$. NN-A003 now gives
(A017.1). Since void probabilities decrease with interval length,
the alternating remainder after $m$ terms has absolute value at
most $E(m+1)$ and the usual alternating sign.

Here are full error bounds for the supplied numerical certificate.
Center the spatial interval and write $x=sa$, with
$a,\xi\in[-1/2,1/2]$. The Fourier-feature map has kernel
\[
 T(a,\xi)=\sqrt s\,e^{2\pi is a\xi},\qquad K_s\simeq TT^*.
\]
Replace the exponential by its Taylor polynomial of degree $p$,
defining $T_p$. With $h=\pi s/2$ and
$R=e^h h^{p+1}/(p+1)!$, the uniform remainder gives
\[
 \|T-T_p\|_{\rm HS}\le\sqrt s R,\quad
 \|T\|_{\rm HS}=\sqrt s,\quad
 \|T_p\|_{\rm HS}\le\sqrt s(1+R).
\]
Thus, for $A=TT^*$ and $B=T_pT_p^*$,
\[
 \|A-B\|_1\le sR(2+R)=:\delta,
 \qquad \operatorname{Tr}B\le s(1+R)^2.               \tag{A017.4}
\]
For positive finite matrices $A,B$, differentiation along their
positive linear interpolation proves
\[
 |\det(I-A)-\det(I-B)|
 \le\|A-B\|_1e^{\max(\operatorname{Tr}A,
                                    \operatorname{Tr}B)}.
\]
To see this even when the interpolation passes through eigenvalue
one, use the polynomial adjugate derivative. If its positive
eigenvalues are $\lambda_i$, the adjugate operator norm is bounded
by $\prod_i(1+\lambda_i)\le e^{\sum_i\lambda_i}$, since
$|1-\lambda_i|\le1+\lambda_i$. Pairing with $A-B$ uses its
trace norm. For positive trace-class operators, compress by common
finite-rank projections and pass to the limit. The determinant
convergence follows from its absolutely convergent exterior-power
series, bounded by $\sum_k\|A\|_1^k/k!$; trace-norm convergence
of the compressions follows first for finite-rank approximants and
then by approximation. Consequently the certificate's more generous
error bound
\[
 \delta\exp(1+s+s(1+R)^2)                             \tag{A017.5}
\]
is valid.

The finite determinant is evaluated without quadrature. Let
\[
 G_{kl}=\int_{-1/2}^{1/2}u^{k+l}du
 =\begin{cases}
  [2^{k+l}(k+l+1)]^{-1},&k+l\ \text{even},\\
 0,&k+l\ \text{odd},
 \end{cases}
 \quad D_{kk}=\frac{(2\pi is)^k}{k!},\quad0\le k,l\le p.
\]
Sylvester's finite-rank determinant identity gives
\[
 \det(I-B)=\det(I-sDGD^*G).
\]
The even and odd index blocks are real and separate. In those
blocks $(DGD^*)_{kl}=(-1)^{(k-l)/2}
 (2\pi s)^k(2\pi s)^lG_{kl}/(k!l!)$.
The program uses exactly these rational-moment matrices with Arb
balls for $\pi$, and enlarges each determinant ball by (A017.5).
For $p=100$ it certifies the six integer gap probabilities, beginning
with
\begin{align*}
 E(1)&=0.17021742137918523073265305256\ldots,\\
 E(2)&=0.00349732514916909769841342811\ldots.
\end{align*}
and ending with
\[
 E(5)=1.55338899144701217136\ldots\,10^{-14},\qquad
 E(6)=1.89553739144083908383\ldots\,10^{-20}.
\]
Their complete intervals and analytic approximation errors are in
\path{cue_agreement_certificate.json}. The odd partial sum
through five gives a lower bound for $A_{\rm adj}$, and the even
partial sum through six an upper bound. Outward rounding gives
exactly (A017.2). The included program asserts both displayed
decimal inequalities as ball comparisons.
\end{proof}

The determinant is the exterior-algebra trace of the sine projection
restricted to an interval. The finite certificate uses an actual
polynomial feature space and its two Gram matrices, not a finite
field or a claimed arithmetic motive. This supplies a rigorously
computable model statistic, with an analytic approximation error;
it does not certify that arithmetic labels have this limit.
